\section{Large entrance scales for the selected service}
\label{ex-centered-geyser-scale-estimates}

This section proves the finite-volume scale estimates used by the selected multiscale service. It isolates the elementary large-parameter estimates and then explains how their separate eventual thresholds are made into one monotone delay sequence. All scale, profile, and gauge data below are the actual fixed data of the selected run.

Write
\[
 B(x)=\left\lceil(\log x)^2\right\rceil,\qquad
 h(x)=e^{-\frac12\log x},\qquad
 f_d(x)=x^{-\left(\deg F_d+4\right)},\qquad
 s(x)=\tfrac14h(x),
\]
and, for $x>1$ and an integer radius $B$,
\[
 \operatorname{stain}(x)=\exp\!\left[-\exp\!\left(\tfrac15(\log\log x)^2\right)\right],
 \qquad \operatorname{siteFactor}(B)=e\{2(16B+9)\}^4.
\]
Here $F_d$ is the finite frame at grade $d$. Fix a nonnegative integer occupancy schedule $Q$ and the positive gauge slope constant $c_G$. For a fixed stage $j$, write $L_j$ for its block side, $b_j=c_G\log L_j$, $y_j(x)=x+b_j$, and $p_j=Q(j)$ for its occupancy allowance. At stage $j$ fix the column count $C_j$ and support exponent $E_j$. For each grade fix its finite local-section constant $A_d$. These quantities remain fixed while $x\to\infty$. The stage index is not a finite-volume period index.

\begin{lemma}[Uniform eventual scale estimates for one transition]
\label{lem:ex-geyser-radial}
Fix stages $i,j$ satisfying $\mathrm{Step}(j,i)$, meaning $i\leq j\leq i+1$, and a finite grade range $0\leq d\leq i$. Let $C_i,C_j\geq1$ be the stage column counts and $E_i,E_j\geq0$ the support exponents. For each grade let $A_d\geq1$ be its uniform local-section constant, $a_d\geq1$ its motion orbit cost, and $m_d=\deg F_d$ its frame degree. Let $L_j\geq1$ be the integer block side and assume the finite title cost $T_j\geq0$. Put
\[
 r_j=2e^{E_j}\,T_j.
\]
Then there is $x_0>1$ so that every $x\geq x_0$ satisfies
\begin{align}
 &8\leq B(x),\qquad y_j(x)^3\leq e^{B(x)/128},\label{eq:ex-geyser-first}\\
 & (2B(x)+1)^4\left(\operatorname{stain}(x)+p_j\,\operatorname{siteFactor}(B(x))s(x)\right)\leq3,\notag\\
 &e^{\operatorname{stain}(x)+p_j\operatorname{siteFactor}(B(x))s(x)}-1\leq1,\label{eq:ex-geyser-second}
\end{align}
and, for every $d\leq i$,
\begin{align}
 &f_d(x)C_i\leq1,\label{eq:ex-geyser-creek1}\\
 &2a_d\bigl(r_j+2C_j\bigr)f_d(y_j(x))\leq1,\label{eq:ex-geyser-creek2}\\
 &e^{E_i+2}A_d x^{m_d}f_d(x)\leq s(x),\label{eq:ex-geyser-creek3}\\
 &y_j(x)^{-1}\leq1,\qquad
 y_j(x)^{-1}L_j^d C_j f_d(x)\leq f_d(y_j(x)),\label{eq:ex-geyser-creek4}
\end{align}
as well as $\nu_d\leq B(x)$, where $\nu_d$ is the minimum stencil side at grade $d$.
\end{lemma}

\begin{proof}
We check each bound from its large-$x$ mechanism. Since $B(x)\to\infty$, eventually $B(x)\geq8$ and $B(x)\geq\nu_d$ for all the finitely many grades. The first polynomial estimate follows from the same elementary comparison: for any fixed $b$ and integer $q$, eventually $x+b\leq2x$, $C2^q\leq x$, and $(q+1)\log x\leq(\log x)^2/128$. Thus
\[
 C(x+b)^q\leq C2^q x^q\leq x^{q+1}
 =e^{(q+1)\log x}\leq e^{B(x)/128}.
\]
Apply this with $C=1,q=3,b=b_j$ to get \eqref{eq:ex-geyser-first}.

For \eqref{eq:ex-geyser-second}, set $H=h(x)$. We verify the reserves
directly. The hat is $H=x^{-1/2}$, so $H\leq1/4$ eventually. Define
\[
 \operatorname{smallK}(x):=\exp\!\left(-\tfrac12
       \exp\!\left(\tfrac15(\log\log x)^2\right)\right).
\]
Writing $y=\log\log x$, for $y\geq10$ we have
$y^2/5-y\geq10>\log6$, hence
$\exp(y^2/5)\geq6e^y=6\log x$. Comparing exponents yields
$\operatorname{smallK}(x)\leq H^6$.

For the contact bound, put $t=\frac12\log x$, so $H=e^{-t}$. For $t\geq1$,
$B(x)=\lceil4t^2\rceil\leq5t^2$, and therefore
\[
2B(x)+1\leq11t^2,\qquad
\operatorname{siteFactor}(B(x))
=e\bigl(2(16B(x)+9)\bigr)^4\leq e(178t^2)^4.
\]
Thus the contact expression is at most $C_p t^{16}e^{-t}$ for a constant
$C_p$ depending only on the fixed occupancy allowance $p_j$. For any fixed integer $q\geq0$,
the exponential series gives $e^t\geq t^{q+1}/(q+1)!$, and so
$C t^qe^{-t}\leq C(q+1)!/t\to0$. Taking $q=16$ proves the contact
reserve. Consequently, for large $x$,
\[
 0<H\leq\tfrac14,\qquad \operatorname{smallK}(x)\leq H^6,
 \qquad (1+p_j\operatorname{siteFactor}(B(x)))(2B(x)+1)^4H\leq1.
\]
The stain is $\operatorname{smallK}(x)^2$. Squaring the small-$K$ bound gives $\operatorname{stain}(x)\leq H^{12}\leq H$ because $0<H\leq1$. Also $s(x)=H/4\leq H$. Therefore
\[
 \operatorname{stain}(x)+p_j\operatorname{siteFactor}(B(x))s(x)
 \leq (1+p_j\operatorname{siteFactor}(B(x)))H,
\]
which proves the product bound by the third reserve estimate. Since $(2B+1)^4\geq2$ once $B\geq8$, the exponent in that display is at most $1/2$. It follows from $1/2<\log2$ that its exponential minus one is at most $e^{1/2}-1<1$.

For the remaining five estimates, fix $d\leq i$ and put $m=\deg F_d+4$, so $m\geq1$, and $D=L_j^d C_j$. All constants are fixed, while $x\to\infty$ and $y=x+b_j\to\infty$. Eventually
\[
 x\geq\max\{1,C_i,4e^{E_i+2}A_d\},\quad
 y\geq\max\{1,2a_d(r_j+2C_j),D2^m\},\quad y\leq2x.
\]
Now $f_d(x)=x^{-m}$, so $C_ix^{-m}\leq1$ and
\[
 2a_d(r_j+2C_j)y^{-m}\leq1.
\]
For the local-section term, put $m=m_d+4=\deg F_d+4$; then
\[
 e^{E_i+2}A_d x^{\deg F_d}f_d(x)
 =\frac{e^{E_i+2}A_d}{x^4}
 \leq\frac{1}{4\sqrt{x}}=s(x),
\]
using $4e^{E_i+2}A_d\leq x\leq x^{7/2}$ for $x\geq1$ (the weaker power inequality suffices). Also $y^{-1}\leq1$.

Finally, $D2^m\leq y$ and $y\leq2x$ imply
\[
 D y^m\leq yx^m.
\]
Indeed $y^m\leq(2x)^m$ and $D2^m\leq y$. Dividing by the positive quantity $y^{m+1}x^m$ yields
\[
 y^{-1}D x^{-m}\leq y^{-m},
\]
which is the last inequality. There are finitely many $d\leq i$, so the finitely many thresholds just chosen have a common maximum. This proves the lemma.
\end{proof}

\begin{theorem}[One monotone delay for all selected transitions]
\label{thm:ex-selected-geyser-steam}
Fix the selected schedule, profile, gauge, stage block sides, and an initial delay sequence $z_j$. There is a coordinatewise larger delay sequence $v_j$, nondecreasing in $j$, such that for every later delay sequence $u_j\geq v_j$, every $j,i,x$ with $\mathrm{Step}(j,i)$ and $u_j\leq x$ satisfies \eqref{eq:ex-geyser-first}--\eqref{eq:ex-geyser-creek4} for the actual profile, as well as the frame-fit inequality $\nu_d\leq B(x)$ for every $d\leq i$.
\end{theorem}

\begin{proof}
For a fixed pair $(j,i)$ with $\mathrm{Step}(j,i)$, apply Lemma~\ref{lem:ex-geyser-radial} to its finite range $d\leq i$. It provides a threshold $t_{j,i}$ above which all its conclusions hold. A fixed $j$ has only two possible stages, $i=j$ and (when $j>0$) $i=j-1$, because the source relation is $i\leq j\leq i+1$. Set $t_j$ to the maximum of the available $t_{j,i}$, and enlarge it to include $z_j$. Define
\[
 v_j=\max_{0\leq \ell\leq j}t_\ell.
\]
Then $v_j$ is nondecreasing and $z_j\leq v_j$. Whenever $v_j\leq x$, we have $t_j\leq x$, so the required estimates hold for either possible stage. If $u\geq v$ coordinatewise and $u_j\leq x$, then $v_j\leq x$, giving the same conclusions. This is the finite-prefix maximum construction used to collect the eventual estimates; no uniformity over unbounded profile ranges is presumed at a fixed scale.
\end{proof}

\paragraph{Use in the service construction.} The theorem gives the exact quantitative service inequalities used by the selected finite producer: the cubic exponential bound, damping and support reserve, five interpolation bounds, and finite stencil fit. In the source these are combined with the selected history, opening, menu, and finite-transfer inputs in the service construction. This section proves the large-entrance estimates and their common monotone threshold; it does not prove the separate selected nonlinear menu response or the paid purchase that supplies the fixed profile data.

\begin{proposition}[From scale estimates to service inequalities]
\label{prop:ex-geyser-service-fields}
Let $x_m$ be the selected hatted scale, $L_m$ its block side, and $B_m=B(x_m)$. Let $E_m$ be its support exponent and $T_m$ its finite title cost. Suppose the selected schedule and delay $u$ satisfy
\[
 x_{m+1}=x_m+b_{J(m+1)},\qquad x_m\geq1,
 \qquad B_{m+1}^2\leq2B_m^2,\qquad L_m\geq2,\qquad u_{J(m+1)}\leq x_m,
\]
with $u_j\geq v_j$ as in Theorem~\ref{thm:ex-selected-geyser-steam}. Assume $\mathrm{threshold}_m>0$ and $\mathrm{threshold}_m^{-1}\leq x_m$. Let $d\leq J(k)$, $k\leq m$, and set
\[
 \zeta_m=f_d(x_m),\quad \phi_m=x_{m+1}^{-1},\quad
 \delta_m=x_{m+1}^{-2},\quad \psi_m=2e^{E_{m+1}}T_{m+1},\quad
 \eta_m=s(x_m).
\]
Write $\Gamma_m$ for the level-$m$ column count, $a_d$ for the finite motion orbit cost, and let $A_d\geq1$ be the same grade-specific local-section constant used in the scale estimates. Then for every $m\geq k$ the scale estimates imply
\begin{gather*}
 \Gamma_m\zeta_m\leq1,\qquad
 2a_d(\psi_m+2\Gamma_{m+1})\zeta_{m+1}\leq1,\qquad \psi_m\geq0,\\
 0<\phi_m\leq1,\quad \delta_m>0,\quad
 \phi_m L_m^d\Gamma_{m+1}\zeta_m\leq\zeta_{m+1},\\
 e^{E_m+2}A_d\,\mathrm{threshold}_m^{-\deg F_d}\zeta_m\leq\eta_m,\\
 1\leq\delta_m\phi_m e^{B_m/128},\quad B_m\geq8,\quad B_{m+1}^2\leq L_mB_m^2.
\end{gather*}
If the local-section bounds are taken from the uniform estimate
$\|\mathrm{places}_d(i,0,U)\|\leq A_d$, its place bound is uniform in $m$;
its footprint cover is at most one once $\nu_d\leq B_m$.
\end{proposition}

\begin{proof}
Apply Theorem~\ref{thm:ex-selected-geyser-steam} at $j=J(m+1)$, $i=J(m)$, and $x=x_m$. The selected clock identities give $\mathrm{Step}(j,i)$ and identify the next scale with $y_j(x_m)=x_{m+1}$. The delay condition places $x_m$ beyond the collected threshold. Since $d\leq J(k)\leq J(m)$, its scale inequalities apply to grade $d$.

Write $F_d$ for the grade frame and $m_d=\deg F_d$. The first scale inequality is
\[
 x_m^{-(m_d+4)}\,\mathrm{columns}_m\leq1,
\]
the first coefficient-budget inequality. Its second is
\[
2\,\mathrm{orbitCost}(F_d)(\psi_m+2\,\mathrm{columns}_{m+1})x_{m+1}^{-(m_d+4)}\leq1,
\]
the second coefficient-budget inequality. The title cost is nonnegative, so $\psi_m\geq0$.

The shifted reciprocal estimate $x_{m+1}^{-1}\leq1$ gives $\phi_m\leq1$. The final scale estimate reads
\[
 x_{m+1}^{-1}L_m^d\mathrm{columns}_{m+1}x_m^{-(m_d+4)}
 \leq x_{m+1}^{-(m_d+4)};
\]
this is the adjacent price inequality with $\phi_m=x_{m+1}^{-1}$ and the adjacent marker prices $\zeta_m,\zeta_{m+1}$.

The local-section bound satisfies $A_d\geq1$. Its weighted norm cost is
\[
 e^{E_m+2}\,A_d\,
      \mathrm{threshold}_m^{-m_d}\,x_m^{-(m_d+4)}.
\]
The threshold comparison and $\mathrm{threshold}_m^{-1}\leq x_m$ bound this by
\[
 e^{E_m+2}\,A_d\,x_m^{m_d}x_m^{-(m_d+4)}
 \leq s(x_m)=\eta_m,
\]
using the third scale inequality. This verifies the weighted norm-cost bound. The separate local-section estimate gives the stated uniform place bound; the footprint estimate follows from the fitted stencil radius and the one-site cover lemma.

Positivity of $\phi_m$ and $\delta_m$ is immediate because $x_{m+1}>0$. The exponential scale inequality follows from the second scale estimate:
\[
 \delta_m\phi_m e^{B_m/128}
 =x_{m+1}^{-3}e^{B_m/128}\geq1.
\]
The contact inequality supplies support $B_m\geq8$. Finally,
\[
 B_{m+1}^2\leq2B_m^2\leq L_mB_m^2,
\]
so the scale split holds. These inequalities supply the scale conditions used by the service certificate.
\end{proof}
